\subsection{Atributos de calidad}
\label{sec:atributos_calidad}

Los atributos de calidad los elegiremos basado en la ISO 25010, que define un modelo de calidad de software. Este modelo se compone de dos partes: características y subcaracterísticas. Las características son las propiedades generales del software, mientras que las subcaracterísticas son propiedades más específicas. Las mismas se muestran en la \reffig{iso_25010}. \\

\defaultfigure{implementacion/iso_25010.png}{Modelo de calidad de software ISO 25010.}{iso_25010}

En este caso, buscamos obtener un software que permita a otros investigadores y desarrolladores trabajar de manera eficiente, intuitiva y directa. Que no requieran un entrenamiento extenso para poder utilizarlo. Por lo tanto, los atributos de calidad que consideraremos son de las características Capacidad de Interaccion, y Adecuación Funcional. Particularmente:

\begin{table}[H]
\centering
\begin{tabular}{|c|c|p{6cm}|}
\hline
\textbf{Requisito} & \textbf{Caracteristica} & \textbf{Descripcion} \\ \hline
Adecuacion Funcional & Pertinencia funcional. & Capacidad del producto software para proporcionar un conjunto de funciones que facilitan la consecución de tareas y objetivos de usuario especificados. \\ \hline
Capacidad de interaccion & Auto-descriptividad. & Capacidad de un producto para presentar la información adecuada, haciendo su uso inmediatamente evidente para el usuario sin interacciones excesivas con el producto u otros recursos (documentación, help desks, otros usuarios...). \\ \hline
\end{tabular}
\caption{Atributos de calidad elegidos}
\label{tab:iso25010}
\end{table}

\subsection{Herramientas existentes}

Una vez entendido el problema a resolver y que atributos de calidad se quieren priorizar, conviene entender cual es el estado del arte hoy en cuanto a herramientas que permitan realizar quantization aware training. Debemos construir una libreria desde cero? Podemos utilizar otras librerias como base? Existe una solucion completa ya disponible que podamos utilizar sin desarrollar ninguna libreria?

Tanto en la industria como en los entornos academicos, existen dos librerias de facto Tensorflow y PyTorch. \\



\subsection{Metodología de trabajo}

\subsubsection{Versionado}
\subsubsection{Estructura del repositorio}
\subsubsection{Testing}
